\section{综合框架与关键结论}

\subsection{硬件演化到软件责任的完整链条}

硬件演化与软件责任构成以下统一框架：
\begin{enumerate}
  \item 摩尔定律长期提高晶体管密度，登纳德缩放长期提高时钟频率；两者共同形成早期单核性能快速增长。
  \item 约在 2005--2006 年，漏电流、功耗与散热约束使频率增长趋于停滞，但晶体管数量仍能继续增加。
  \item 新增晶体管被组织为多核、向量执行资源和专用加速器。现代系统因此具有显著并行性与异构性。
  \item 中央处理器面向低延迟，适合串行依赖和复杂控制；图形处理器面向高吞吐量，适合大量规则且相互独立的计算。
  \item 图形处理器原有的游戏市场和图形并行结构提供了经济与技术基础；统一计算设备架构进一步把它转化为通用计算平台。
  \item 应用必须显式暴露并行性。性能依赖算法可扩展性、线程规则性、负载均衡、数据布局、有效带宽和争用控制。
  \item 任意局部加速最终都受完整应用中未加速部分限制，阿姆达尔定律给出这一上界。
\end{enumerate}

\subsection{选择中央处理器或图形处理器时的判断顺序}

面对一个待加速代码段，可以按下列问题逐步判断：
\begin{enumerate}
  \item 它在完整程序中占多少时间？占比过低时，局部优化的整体收益受限。
  \item 是否存在大量相互独立或弱依赖的任务？任务数量必须足以填充图形处理器。
  \item 各任务是否执行相同或相近的操作？高度不规则控制流会降低并行执行效率。
  \item 每个任务的工作量是否接近？最慢任务会决定并行阶段完成时间。
  \item 数据访问是否规则、具有局部性，并能利用全局内存带宽？
  \item 把数据移动到图形处理器并启动内核的开销，是否小于并行计算所节省的时间？
  \item 实现能否在更大输入和后续硬件上保持正确性、性能与可维护性？
\end{enumerate}

只有当这些问题的总体答案支持加速器映射时，图形处理器才可能产生端到端收益。对于串行依赖明显、任务规模很小或数据移动占主导的代码，中央处理器可能更合适。

\Needspace{10\baselineskip}
\subsection{高频概念辨析}

\begingroup
\small
\begin{longtable}{L{0.34\textwidth}L{0.60\textwidth}}
\toprule
\textbf{容易形成的误解} & \textbf{准确理解} \\
\midrule
\endfirsthead
\toprule
\textbf{容易形成的误解} & \textbf{准确理解} \\
\midrule
\endhead
摩尔定律等同于处理器性能翻倍。 & 历史性能增长来自晶体管密度、频率缩放和处理器结构改进的共同作用；摩尔定律主要描述晶体管密度趋势。 \\
处理器核心数量增加会自动加速串行程序。 & 串行指令流通常只使用一个核心；程序必须暴露并行任务并承担调度、通信和同步成本。 \\
图形处理器的单个线程一定比中央处理器线程快。 & 图形处理器优化大批线程的总吞吐量，单线程延迟可能更高。 \\
统一计算设备架构核心与中央处理器核心可以按数量直接比较。 & 两者粒度不同；统一计算设备架构核心更接近细粒度算术执行通道。 \\
硬件峰值内存带宽就是程序能够获得的带宽。 & 实际有效带宽取决于数据布局、访问模式、复用、并发程度和争用。 \\
并行部分加速 1000 倍会让整个程序加速 1000 倍。 & 整体加速由 $S=1/(1-p+p/s)$ 决定，未加速部分最终成为上界。 \\
问答系统使用课程材料，因此答案必然正确。 & 检索增强生成只能改善相关性；检索遗漏、隐含假设与生成错误仍可能导致错误答案。 \\
\bottomrule
\end{longtable}
\endgroup

\subsection{“高性能”的操作性定义}

高性能需要放在完整软件工程目标中理解。并行实现应满足：
\begin{itemize}[itemsep=0.15em,parsep=0pt,topsep=0.3em]
  \item \textbf{功能正确：}与参考语义一致，能够通过规定的数据集与问题检查；
  \item \textbf{端到端加速：}完整应用时间缩短，而不是只报告某个局部内核的理想时间；
  \item \textbf{可解释：}能够用并行度、访存、负载、串行比例等因素解释性能变化；
  \item \textbf{可扩展：}问题规模或硬件资源增加时仍能利用更多并行能力；
  \item \textbf{可维护与可移植：}实现结构清晰，不依赖难以复现的偶然参数，并能适应硬件代际变化。
\end{itemize}

这五项要求也解释了项目评分为何同时关注功能、编码风格和性能。单独追求某一项会产生失衡：正确但极慢的实现没有利用加速器价值；快速但错误的实现改变了问题；高度手工特化但不可维护的实现难以跨代延续。

\subsection{阅读与操作任务}

需要完成以下事项：
\begin{itemize}[itemsep=0.15em,parsep=0pt,topsep=0.3em]
  \item 阅读教材第 1 章；
  \item 阅读《统一计算设备架构编程指南》第 1.1 节“导论”；
  \item 尽快创建高等科学与工程研究计算环境账户并取得 Delta 访问权限；
  \item 配置个人代码仓库，完成课程发布代码的同步流程；
  \item 在账户可用后验证安全外壳协议登录、编译节点与计算节点作业提交流程。
\end{itemize}

\begin{keybox}
软件可持续性要求算法与库具备可扩展性和可移植性，从而跨越多代硬件；其价值来自对并行结构和性能瓶颈的清晰建模，而不是依赖某一代设备的偶然峰值。
\end{keybox}
